\documentclass[10pt]{article}
\usepackage[T1]{fontenc}
\usepackage[utf8]{inputenc}
\usepackage[french]{babel}
\usepackage{amsmath,amssymb,mathrsfs,mathpazo}
\usepackage{geometry}
\geometry{a4paper,margin=2cm}
\usepackage{graphicx,array,tabularx,tabularray,colortbl,multirow,multicol,booktabs,enumitem,pifont,ulem,fancybox,diagbox,cancel}
\usepackage{pstricks,pstricks-add,pst-plot,pst-tree,pst-node,pst-text,pst-func,pst-math,pst-bspline,pst-3dplot,pst-eucl,pst-solides3d}
\let\pscalculate\undefined
\usepackage{pst-bezier}
\usepackage{tikz,pgfplots,tkz-tab}
\usetikzlibrary{arrows,arrows.meta,patterns,calc,positioning,trees,decorations.markings,shapes.geometric,shapes.misc,shapes.arrows,intersections,decorations.pathmorphing,decorations.pathreplacing,backgrounds,shadows,babel}
\pgfplotsset{compat=1.18}
\usepgfplotslibrary{fillbetween}
\usepackage[np]{numprint}
\usepackage{listings,makecell,pas-tableur,twoopt}
\let\barre\undefined
\usepackage[pstricks]{tablvar}
\definecolor{darkgreen}{RGB}{0,100,0}
\providecommand{\red}{\color{red}}
\providecommand{\blue}{\color{blue}}
\providecommand{\green}{\color{green}}
\providecommand{\black}{\color{black}}
\providecommand{\white}{\color{white}}
\providecommand{\gray}{\color{gray}}
\providecommand{\cyan}{\color{cyan}}
\providecommand{\magenta}{\color{magenta}}
\providecommand{\yellow}{\color{yellow}}
\providecommand{\violet}{\color{violet}}
\providecommand{\orange}{\color{orange}}
\providecommand{\euro}{\texteuro}
\providecommand{\EUR}{\texteuro}
\providecommand{\ieme}{\textsuperscript{e}}
\providecommand{\degres}{\ensuremath{^\circ}}
\providecommand{\degree}{\ensuremath{^\circ}}
\providecommand{\vv}[1]{\overrightarrow{#1}}
\providecommand{\Coord}[1]{#1}
\providecommand{\decosix}{}
\providecommand{\starredbullet}{\ensuremath{\bullet}}
\pagestyle{empty}
\setlength{\parindent}{0pt}
\usepackage{adjustbox}\DeclareUnicodeCharacter{2236}{:}\DeclareUnicodeCharacter{2190}{\ensuremath{\leftarrow}}
\providecommand{\R}{\mathbb{R}}
\providecommand{\N}{\mathbb{N}}
\providecommand{\Z}{\mathbb{Z}}
\providecommand{\Q}{\mathbb{Q}}
\providecommand{\D}{\mathbb{D}}
\providecommand{\C}{\mathbb{C}}
\providecommand{\e}{{\mathrm{e}}}
\providecommand{\vect}[1]{\overrightarrow{#1}}
\providecommand{\vectt}[1]{\overrightarrow{\mathrm{#1}}}
\providecommand{\vvt}[1]{\overrightarrow{\mathrm{#1}}}
\providecommand{\np}[2][]{{#2~#1}}
\providecommand{\vv}[1]{{\overrightarrow{#1}}}
\renewcommand{\d}{{\mathrm{d}}}
\providecommand{\degres}{{^\circ}}
\providecommand{\degre}{{^\circ}}
\providecommand{\blue}{\color{blue}}
\providecommand{\red}{\color{red}}
\providecommand{\enskip}{\quad}
\providecommand{\llbracket}{\lbrack\!\lbrack}
\providecommand{\rrbracket}{\rbrack\!\rbrack}
\providecommand{\euro}{{\text{€}}}
\providecommand{\no}{{\text{n°}}}
\providecommand{\interval}[2]{{[#1;#2]}}
\providecommand{\intervFO}[2]{{[#1;#2[}}
\providecommand{\dotfill}{\ldots}

\providecommand{\ds}{\displaystyle}
\providecommand{\textding}[1]{#1}
\providecommand{\Oij}{\(\mathrm{O};\vec{\imath},\vec{\jmath}\)}
\providecommand{\Oijk}{\(\mathrm{O};\vec{\imath},\vec{\jmath},\vec{k}\)}
\providecommand{\Ouv}{\(\mathrm{O};\vec{u},\vec{v}\)}
\providecommand{\intervFF}[2]{\ensuremath{[#1;#2]}}
\renewcommand{\euro}{\ifmmode\text{\texteuro}\else\texteuro\fi}
\DeclareUnicodeCharacter{00B0}{\ifmmode^{\circ}\else\textdegree\fi}
\DeclareUnicodeCharacter{20AC}{\ifmmode\text{\texteuro}\else\texteuro\fi}
% Énoncés : Éducation nationale. Transcription et corrigés : APMEP, Denis Vergès et contributeurs. Découpage et classement : CoopMaths, Sébastien Lozano, Éric Elter et contributeurs. Adaptation de la présentation : Maths Explorer.
% CC BY-SA : https://coopmaths.fr/www/about/
% https://coopmaths.fr/alea/static/e3c/2020/tex/e3c_2020_00_sujet39_2_cor.tex
\begin{document}
\noindent\textit{APMEP / CoopMaths. CC BY-SA. Présentation adaptée par Maths Explorer.}\par\medskip
\subsection*{1.}

La fonction \( f(x) = x^2 - 3x + 4 \) est un polynôme dérivable sur \( \mathbb{R} \), donc sur \( [0\,;\,+\infty[ \). Sur cet intervalle :
\[
f'(x) = 2x - 3.
\]

\(2x - 3 < 0 \iff x < \dfrac{3}{2}\), donc \(f\) est décroissante sur \( \left]0\,;\,\dfrac{3}{2} \right[\) puis croissante sur \( \left]\dfrac{3}{2}\,;\,+\infty \right[ \).

Le minimum est :
\[
f\left(\dfrac{3}{2}\right) = \dfrac{9}{4} - \dfrac{9}{2} + 4 = \dfrac{7}{4}.
\]

\subsection*{2.}

\paragraph{a.} \(\mathcal{C}\) est la représentation graphique de la fonction \(x \longmapsto \sqrt{x}\), donc \( M(x \,;\, \sqrt{x}) \).

\paragraph{b.} On a donc :
\begin{align*}
AM^2 &= (x - 2)^2 + (\sqrt{x} - 0)^2 \\
&= x^2 + 4 - 4x + x \\
&= x^2 - 3x + 4.
\end{align*}

\paragraph{c.} Pour le trinôme \(x^2 - 3x + 4\), on a :
\[
\Delta = 9 - 16 = -7 < 0.
\]
Ce trinôme n'a pas de racines et on sait que son minimum est atteint pour :
\[
x = -\dfrac{b}{2a} = -\dfrac{-3}{2} = \dfrac{3}{2}.
\]
Donc le point correspondant au point de \(C\) le plus proche de \(A\) a pour coordonnées \( \left( \dfrac{3}{2} \,;\, \sqrt{\dfrac{3}{2}} \right).  \textit{ Ce point est noté } B \textit{ pour la suite}\).

\paragraph{d.} On a pour \(x \neq 0\), \(f'(x) = \dfrac{1}{2\sqrt{x}}\).

La tangente en \(B\) a pour coefficient directeur :
\[
f'\left( \dfrac{3}{2} \right) = \dfrac{1}{2 \sqrt{\dfrac{3}{2}}} = \dfrac{1}{\sqrt{2} \times \sqrt{3}} = \dfrac{1}{\sqrt{6}}.
\]
La droite \((AB)\) a pour coefficient directeur :
\[
\dfrac{\sqrt{\dfrac{3}{2}} - 0}{\dfrac{3}{2} - 2} = \dfrac{\sqrt{\dfrac{3}{2}}}{-\dfrac{1}{2}} = -2 \sqrt{\dfrac{3}{2}}.
\]
Le produit des deux coefficients directeurs est :
\[
\dfrac{1}{\sqrt{6}} \times \left( -2 \sqrt{\dfrac{3}{2}} \right) = \dfrac{-2\sqrt{3}}{\sqrt{2} \times \sqrt{3} \times \sqrt{2}} = \dfrac{-2}{2} = -1.
\]
L'élève a raison, les droites sont bien perpendiculaires.

\end{document}
